\subsection{DET-19: the arithmetic flow is not the natural Bessel boost flow}
The realization of \(B\) on the source Bessel space does not identify
its arithmetic translation group with the source group's natural
noncompact one-parameter action. The actual decay hypothesis gives
a direct obstruction to that stronger correspondence. It applies
already to one explicit subgroup, without an assumption of full
representation invariance.

Retain the \texttt{Std} data of \fid{DET-15} on \(H_\sigma\), and
write \(D=\operatorname{span}\{\sigma(g)w_0:g\in G\}\), where
\(G=\mathrm{UH}(H_{k,r})\), \(k\ge1\), and \(n=2r+k+1\).
The source
\path{Development/AbelianVarietyTheory/Definitions/SemilocalFrames/MatrixOfBesselTwoSidedAction.lean},
line 15, defines \(\operatorname{opn}(g)\) as its operator norm on
Euclidean \(\C^n\); lines 50--55 give
\begin{equation}\label{eq:boost-source-decay}
 a\ge2n^5>0,\qquad
 |\langle v,\sigma(g)u\rangle|
 \le C_{u,v}\operatorname{opn}(g)^{-a},
 \quad u,v\in D,\quad g\in G.
\end{equation}
The constants may depend on \(u,v\). No uniform decay rate on
the whole Hilbert space is presumed. The source unitary and strong
continuity conditions are
\path{Formalism/SharpRestriction/CuspSpanNormBounds/RelevantZeroGramTerms.lean},
lines 23--28.

Define \(g_t\) to be the identity outside the source middle
coordinates \(r,r+k\), and on those two coordinates put
\[
 g_t=\begin{pmatrix}\cosh t&\sinh t\\
                      \sinh t&\cosh t\end{pmatrix},\qquad
 V_t=\sigma(g_t),\qquad t\in\R.
\]
The two coordinates are distinct because \(k\ge1\).

\begin{proposition}
Under the retained source hypotheses:
\begin{enumerate}
\item \(V_t\) is a strongly continuous unitary group and
\[
 \langle v,V_tu\rangle\longrightarrow0
 \quad(|t|\longrightarrow\infty)
 \qquad(u,v\in H_\sigma).
\]
For every trace-class \(T\) on \(H_\sigma\), the continuous
function \(t\mapsto\Tr(TV_t)\) also tends to zero.
\item There is no nonzero boost eigenvector
\(V_tu=e^{i\lambda t}u\), no nonzero finite-dimensional invariant
subspace, and no nonzero compact operator commuting with every
\(V_t\).
\item The arithmetic kernel \(k\) of \fid{DET-16} does not tend
to zero at infinity, unconditionally. Consequently it is neither
\(\langle v,V_tu\rangle\) nor \(\Tr(TV_t)\) for any indicated
vectors or trace-class \(T\). The same conclusion holds for
\(2\Tr(T(V_t+V_{-t})/2)\).
\item Under RH, every bounded intertwiner from the cyclic
arithmetic translation model of \fid{DET-16}, or its restored
full-multiplicity model, into \((H_\sigma,V_t)\) is zero.
\end{enumerate}
All these conclusions transfer to the isometric Bessel submodel
of \(H_\sigma\). Only that submodel's right-translation action
is asserted here.
\end{proposition}
\begin{proof}
The displayed definition of \(H_{k,r}\) in
\path{Formalism/FormalDegreeBessel/StatementG.lean}, lines 20--24,
makes the indicated two-plane orthogonal to the other coordinates
and gives it matrix \(\operatorname{diag}(1,-1)\).
The identity \(\cosh^2t-\sinh^2t=1\) therefore gives
\(g_t^*H_{k,r}g_t=H_{k,r}\). Thus \(g_t\in G\), and the
hyperbolic addition formulas give \(g_tg_s=g_{t+s}\).
This is a continuous group map. The source representation
conditions make \(V_t\) a strongly continuous unitary group.
On the Euclidean orthonormal vectors
\((e_r+e_{r+k})/\sqrt2\) and
\((e_r-e_{r+k})/\sqrt2\), the matrix \(g_t\) has eigenvalues
\(e^t,e^{-t}\); it is the identity on their orthogonal
complement. Hence
\[
 \operatorname{opn}(g_t)=e^{|t|}.
\]
For \(u,v\in D\), \eqref{eq:boost-source-decay} now bounds the
coefficient by \(C_{u,v}e^{-a|t|}\), proving its vanishing.

The space \(D\) is dense. Its closure is a nonzero closed
invariant subspace, because \(w_0\in D\), \(\|w_0\|=1\),
and every \(\sigma(g)\) is bounded and preserves the orbit
span. Source irreducibility makes the closure \(H_\sigma\).
For arbitrary \(u,v\in H_\sigma\), choose \(u_0,v_0\in D\).
Unitarity gives the uniform estimate
\[
 |\langle v,V_tu\rangle-\langle v_0,V_tu_0\rangle|
 \le\|v-v_0\|\|u\|+\|v_0\|\|u-u_0\|.
\]
Make its right side arbitrarily small, then let \(|t|\) tend
to infinity in the coefficient on \(D\). This proves vanishing
for all vectors, without claiming their exponential rates.

If \(R\) is finite rank, write
\(Rx=\sum_{j=1}^q\langle v_j,x\rangle u_j\).
The conjugate-linear-first convention gives
\[
 \Tr(RV_t)=\sum_{j=1}^q\langle v_j,V_tu_j\rangle,
\]
so this continuous function vanishes at infinity.
For \(T\in\mathcal S_1\), choose finite-rank \(R\) with
\(\|T-R\|_1<\epsilon\). Uniformly in \(t\),
\[
 |\Tr((T-R)V_t)|\le\|T-R\|_1\|V_t\|<\epsilon.
\]
The uniform limit is continuous and vanishes at infinity.
This proves the first assertion.

A nonzero eigenvector would give
\(|\langle u,V_tu\rangle|=\|u\|^2>0\), contradicting the
coefficient limit. More generally, if a nonzero finite-dimensional
space \(E\) were invariant under every \(V_t\), invariance under
\(V_{-t}\) would make \(V_t|_E\) unitary. For an orthonormal
basis \(e_1,\ldots,e_m\) of \(E\),
\[
 \sum_{i,j=1}^m|\langle e_i,V_te_j\rangle|^2=m.
\]
The left side tends to zero by the first assertion, a
contradiction. If a nonzero compact \(Q\) commuted with all
\(V_t\), its adjoint would also commute, since \(V_t^*=V_{-t}\).
Then \(Q^*Q\) would be nonzero positive compact and commute
with the group. Its largest positive eigenvalue has a nonzero
finite-dimensional eigenspace, invariant under every \(V_t\).
This is impossible, proving the compact-commutant assertion.

For the arithmetic assertion, a bounded continuous function
\(h(t)\) tending to zero at infinity has zero Cesaro coefficient
at every real frequency: for any \(L>0\),
\[
 \left|\frac1{2T}\int_{-T}^{T}h(t)e^{-i\lambda t}\,dt\right|
 \le\frac{1}{2T}\int_{-L}^{L}|h(t)|\,dt
      +\sup_{|t|>L}|h(t)|\qquad(T>L).
\]
First let \(T\to\infty\), then \(L\to\infty\).
Suppose \(k\) tended to zero. It would be bounded, so
\fid{DET-16}, using the complete global Landau proof of
\fid{STEP6-03}, would imply RH. Under RH, the exact explicit
formula and Cesaro calculation in \fid{DET-16} give, at every
positive Xi ordinate \(\lambda\),
\[
 \lim_{T\to\infty}\frac1{2T}
       \int_{-T}^{T}k(t)e^{-i\lambda t}\,dt
 =\frac{m_\lambda}{\lambda^2+1/4}>0.
\]
There are such zeros by the classical zero count in \fid{B-01}.
This contradicts the preceding estimate. Thus \(k\) does not
tend to zero, without an RH assumption in this conclusion.
All the proposed source trace coefficients, including their
displayed cosine combination, do tend to zero and therefore
cannot equal \(k\).

Under RH, the atomic model in \fid{DET-16} has a complete
orthonormal basis \(e_\lambda\) satisfying
\(U_te_\lambda=e^{i\lambda t}e_\lambda\); in the restored
model there are exactly \(m_K(\lambda)\) such basis vectors
at each retained frequency. If a bounded map \(J\) satisfies
\(JU_t=V_tJ\), every \(Je_\lambda\) is a boost eigenvector,
hence zero. Density of the basis and boundedness give \(J=0\).
This excludes, in particular, an injective bounded intertwiner
or a unitary identification of these flows.

Finally, let
\(\mathcal J:H_\sigma\longrightarrow\mathcal B_\sigma
\subset\mathcal B\)
be the isometry from \fid{DET-15}, defined on \(D\) by
\(\mathcal Ju=\eta_u=\sqrt d\,[x\mapsto l(\sigma(x)u)]_\nu\).
Its image \(\mathcal B_\sigma\) is closed. For \(u\in D\)
and \(g,x\in G\), the source representation law gives
\[
 \eta_u(xg)=\sqrt d\,l(\sigma(xg)u)
          =\sqrt d\,l(\sigma(x)\sigma(g)u)
          =\eta_{\sigma(g)u}(x).
\]
This identity concerns the canonical coefficient representatives
on \(G\), which supports \(\nu\). The source orbit span is
invariant, so its right side is again a coefficient of a vector
in \(D\). It defines an isometric right-translation action on
these classes and extends by the isometry to
\(\mathcal J\sigma(g)\mathcal J^{-1}\) on
\(\mathcal B_\sigma\).
Consequently the boost action there is unitarily \(V_t\), and
all the preceding conclusions transfer. We use no assertion
that right translation is unitary or even defined on arbitrary
classes of the ambient \(L^2(\nu;W_p)\).
\end{proof}

In particular, the nonzero positive trace-class operator obtained
in \fid{DET-16} on this Bessel image cannot commute with the
natural boost. Its determinant remains \(G\), and its finite
source convolution approximants remain valid. The obstruction
concerns a proposed identification of flows, not those operator
and determinant constructions. It leaves open a separately
specified historical matrix correspondence and other choices of
flow; it is not an impossibility theorem for arithmetic
realizations of \(B\).
